% !TEX root = ../main.tex
%%%%%%%%%%%%%%%%%%%%%%%%%%%%%%%%%%%%%%%%%%%%%%%%%%%%%%%%%%%%%%%%%%%%%%%
% Repairing the quantity of interest
%%%%%%%%%%%%%%%%%%%%%%%%%%%%%%%%%%%%%%%%%%%%%%%%%%%%%%%%%%%%%%%%%%%%%%%

\section{Repairing the quantity of interest}
\label{sec:bridge}

The contradiction of \cref{sec:inversion} is resolvable only by experiment,
and the experiment is designed as a $2\times2$: method (chord versus OLS)
crossed with conditioning (all other parameters fixed at their defaults versus
marginalised over the sensitivity ranges), on a shared grid of runs. The
verdict rule was fixed \emph{before} execution: a sign flip across the method
axis at fixed parameters implicates cause~(a); a flip across the conditioning
axis at constant method implicates cause~(b); both flipping means both
contribute and the weights must be quantified; neither flipping is a new
finding to be reported, not forced.

\subsection{The fixed-parameter arm}

\Cref{tab:bridge} reports the arm that has been executed.

\begin{table}[htbp]
\centering
\small
\caption{Bridging experiment, fixed-parameter arm. The canonical support is
the committed one; \emph{grid4} is the four-point grid shared with the
marginalised arm.}
\label{tab:bridge}
\begin{tabular}{llrl}
\toprule
Scenario & Method (support) & $\sigstar$ & 95\% CI\\
\midrule
\multirow{4}{*}{$\eta = 0$}
  & OLS (canonical)   & 0.654 & $[0.616,\,0.691]$\\
  & chord (canonical) & 0.447 & $[0.395,\,0.501]$\\
  & OLS (grid4)       & 0.643 & $[0.596,\,0.701]$\\
  & chord (grid4)     & 0.447 & $[0.395,\,0.501]$\\
\addlinespace[2pt]
\multirow{2}{*}{$\eta = 0.10$}
  & OLS (canonical)   & 0.725 & $[0.697,\,0.745]$\\
  & chord (canonical) & 0.501 & $[0.465,\,0.553]$\\
\bottomrule
\end{tabular}
\end{table}

The mandatory anchoring control passes: the OLS estimator on the canonical
support reproduces $\sigstar = 0.654$ with the published interval, so the
problem is not upstream of both candidate causes.

Cause~(a) is then confirmed as a real and \emph{large} contributor. At
identical parameters, changing only the estimator moves $\sigstar$ from
$0.654$ to $0.447$, with non-overlapping intervals: for any $\sigma$ between
those values, the chord is negative while the OLS slope is positive. The two
statistics genuinely disagree in sign on the same data.

\paragraph{But cause (a) points the wrong way.} This is the informative part.
The chord turns negative at a \emph{lower} $\sigma$ than the OLS slope, so at
fixed parameters the chord finds wage-led outcomes \emph{more} often at low
$\sigma$, not less. \Cref{tab:widesigma}, by contrast, finds wage-led outcomes
\emph{rare} below $\sigma \approx 0.65$. The chord/derivative distinction is
therefore real but insufficient: on its own it would have widened the wage-led
region at low $\sigma$, whereas the marginalised measurement narrows it.
Cause~(b) --- conditioning --- must be doing the opposing work.

\begin{declbox}[frametitle={Status: open, and reported as open}]
The marginalised arm of the design has not yet been executed, so the formal
verdict is not yet available and we do not state one. What has been measured
is that (a) contributes materially and in the \emph{opposite} direction to
the observed inversion, which raises the prior on (b) considerably but does
not establish it. Until the marginalised arm exists, the correct description
of $\sigstar$ is a \emph{conditional} result whose marginal counterpart is
unmeasured --- and that is how it is described in this paper. This is the
point at which the project is most exposed to asserting the opposite of the
truth, and it is deliberately left unasserted.
\end{declbox}

\subsection{A quantity of interest suited to a U-shaped response}

The deeper repair is to stop summarising a U-shaped response with a single
slope. We therefore report the location of the turning point $\rhostar$ ---
the argument minimum of the quadratic fit --- alongside the slope, and, most
usefully, the \emph{side} on which the empirically anchored range of $\rho$
falls (\cref{tab:turning}).

\begin{table}[htbp]
\centering
\small
\caption{Turning point of $Y(\rho)$ and the position of the empirically
anchored range $\rho \in [0.35, 0.45]$ relative to it, at $\eta = 0$.}
\label{tab:turning}
\begin{tabular}{lrrrl}
\toprule
$\sigma$ & OLS slope & Chord & $\rhostar$ & Empirical $\rho$ lies\\
\midrule
0.30 & $+41.9$ & $+14.6$ & 0.420 & straddling\\
0.40 & $+26.8$ & $+5.2$  & 0.436 & straddling\\
0.50 & $+17.8$ & $-5.9$  & 0.462 & left of the turn\\
0.60 & $+5.6$  & $-20.7$ & 0.489 & left of the turn\\
0.70 & $-4.8$  & $-18.9$ & 0.513 & left of the turn\\
0.80 & $-11.7$ & $-33.3$ & 0.527 & left of the turn\\
1.00 & $-29.6$ & $-51.5$ & 0.569 & left of the turn\\
1.25 & $-45.3$ & $-68.8$ & 0.592 & left of the turn\\
1.50 & $-61.1$ & $-83.9$ & 0.640 & left of the turn\\
\bottomrule
\end{tabular}
\end{table}

\Cref{tab:turning} reconciles the two statistics mechanically. At
$\sigma = 0.5$ the turning point sits at $\rhostar = 0.462$, inside the
support: the OLS slope over $[0.35, 0.65]$ is positive because the rising
right branch is longer than the falling left branch, while the chord over
$[0.35, 0.55]$ is negative because it is weighted toward the left branch.
Neither is wrong; they are answers to different questions.

The economically interpretable statement is the third column. For every
$\sigma \geq 0.5$, the empirically anchored range of $\rho$ lies \emph{left of
the turning point} --- on the falling branch. The defensible formulation is
therefore not ``the economy is wage-led'' or ``profit-led'' but:

\begin{declbox}
\emph{Retention depresses output while $\rho$ remains below a
$\sigma$-dependent threshold $\rhostar$, and expands it beyond;
$\rhostar$ rises with $\sigma$, from $0.46$ at $\sigma = 0.5$ to $0.64$ at
$\sigma = 1.5$; and the empirically anchored retention rate lies below
$\rhostar$ for every $\sigma$ at or above $0.5$.}
\end{declbox}

This statement does not depend on which chord an analyst happens to pick, which
is exactly the property the original quantity of interest lacked.
